\documentclass[../../../main.tex]{subfiles}

\begin{document}

\section{Additional Topics}
Involution, monotonicity


Some interesting problems gives bound for the function for us to
use bounding argument.

\begin{important}
The key to these kind of problems is to perform certain operations on
the function and show that it contradicts with the given criteria.
\end{important}

Here is an interesting problem from the document
\textit{169 Functional Equations} by Moubinool the showcase the
bounding through a specific criteria given by the problem.

\begin{problem}{}{}
Find all functions $f: \mathbb{R} \to \mathbb{R}$ such that
$\left\{ \frac{f(x)}{x} \mid x \in \mathbb{R} \text{ and } x \neq 0
\right\}$ is a finite set and $f(x - 1 - f(x)) = f(x) - x - 1$ for all
real $x$.
\end{problem}
\begin{solution}
For this problem, because the terms inside $f$ in the left-hand
side and the terms in the right-hand side are similar, let's try
and make them equal like the following.
\[
f(x - 1 - f(x)) = -(x - 1 - f(x)) - 2
\]
Let $\alpha = x - 1 - f(x)$ and we get $f(\alpha) = -\alpha - 2$.
But there is a problem here. The number of terms for $\frac{f(x)}{x}$
is bounded, but we have infinitely many $\frac{-\alpha - 2}{\alpha}$.
This tells us that $\alpha$ is finite for the condition to be
satisfied.

Further observing the problem, let's substitute $\alpha$ for $x$.
We obtain
\begin{align*}
    f(\alpha - 1 - f(\alpha))
    &= -(\alpha - 1 - f(\alpha)) - 2 \\
    f(\alpha - 1 + \alpha + 2) &= -\alpha - 1 - \alpha - 2
\end{align*}
and $f(2\alpha + 1) = -(2\alpha + 1) - 2$. This shows that
$2\alpha + 1$ is also an element of a finite set, say $A$,
that contains $\alpha$.

Continuing, we can try and substitute $2\alpha + 1$ for $x$.
We obtain $f(2\alpha + 1 - 1 + 2\alpha + 1 + 2) = -(2\alpha + 1) - 2
- 2\alpha - 1 - 1$ and $f(4\alpha + 3) = -(4\alpha + 3) - 2$.
This shows that $4\alpha + 3 \in A$. From the patterns, we could
hypothesize that $2^n \alpha + 2^n - 1 \in A$. To show this,
it suffices to show that if $2^{n-1} \alpha + 2^{n-1} - 1 \in A$,
then $2^n \alpha + 2^n - 1 \in A$. Consider the following equations
that we get for substituting $x = 2^{n-1} \alpha + 2^{n-1} - 1$.
\begin{align*}
    &\quad\ f \left(
        2^{n-1} \alpha + 2^{n-1} - 1 - 1
            + \left( 2^{n-1} \alpha + 2^{n-1} - 1) + 2 \right)
    \right) \\
    &= -\left( 2^{n-1} \alpha + 2^{n-1} - 1 \right) - 2
        - \left( 2^{n-1} \alpha + 2^{n-1} - 1 \right) - 1 \\
    &\qquad\quad
    f \left( 2^n \alpha + 2^n - 1 \right)
    = -\left( 2^n \alpha + 2^n - 1 \right) - 2
\end{align*}
Indeed by induction, $2^n \alpha + 2^n - 1 \in A$ for all positive
integers $n$. In other words, $\alpha = -1$ is the only possible
case and $\alpha = x - 1 - f(x) = -1$. Therefore, $f(x) = x$ is the
only solution.
\end{solution}


\end{document}


